\section{Research design and data}\label{sec:data}

\subsection{How the work was organised}
A common risk in empirical work is that the hypotheses, the specification and the data are adjusted to each other after
the results are seen. To limit this, the work followed a fixed sequence (\cref{fig:design}). The model came first. Its
predictions, the samples, the variable definitions, the estimating equations and the results that would count against
each prediction were then written into a pre-analysis plan, which I froze with a time-stamped tag in the project's
version control on 5 November 2026. Until then I had looked only at the data on bans and position limits and at the
basis on its own, never at how one related to the other. Where the plan left a detail open, I wrote the choice into a
dated log before running the test that needed it. Each estimator was first run on data simulated with a known effect, to
check that it recovers that effect. Every pre-specified test is reported, including those that go against the model.

\begin{figure}[t]
\caption{Theory and tests were fixed before any outcome was related to any treatment}\label{fig:design}
\exsub{The sequence of the work, from model to results}
\centering
\begin{tikzpicture}[node distance=4mm,
  box/.style={draw=accent,line width=0.7pt,rounded corners=2pt,text width=2.35cm,minimum height=1.45cm,align=center,font=\footnotesize},
  arr/.style={-{Stealth[length=2.2mm]},line width=0.8pt,color=accent}]
\node[box] (a) {\textbf{Model}\\ every limit a parameter; predictions and falsifiers};
\node[box,right=of a] (b) {\textbf{Frozen plan}\\ samples, tests and signs fixed with a time stamp};
\node[box,right=of b] (c) {\textbf{Data}\\ built from exchange files, checked against the rules};
\node[box,right=of c] (d) {\textbf{Estimators}\\ tested on data simulated with known effects};
\node[box,right=of d,fill=tilebg] (e) {\textbf{Results}\\ all reported, including those against the model};
\foreach \x/\y in {a/b,b/c,c/d,d/e} \draw[arr] (\x) -- (\y);
\end{tikzpicture}
\exsource{The plan and its log are described in \cref{app:plan}.}
\end{figure}

The plan names four primary hypotheses: that costs set the edges of the band (H3), that trading leaves futures in a ban
(H8), that correction becomes one-sided in a ban (H7), and that access to short selling speeds price adjustment (H6).
Their p-values are corrected jointly by Holm's method at the 5\% level. The other tests are secondary and are reported
without adjustment. This report presents H3 and H8 together with the secondary tests that belong to the same questions;
H6 and H7 belong to the second stage.

\subsection{Data}
All data come from public archives of the NSE and are downloaded and processed by scripts, so every step can be repeated
(\cref{tab:sources}). The main files are the daily combined open-interest files, which give each stock's MWPL and open
interest from February 2010; the daily lists of stocks in the ban; and the daily price files (bhavcopies) of the
futures and cash markets from 2015 to October 2026.

\begin{table}[t]
\caption{Every series comes from a public archive}\label{tab:sources}
\exsub{Sources, content and coverage of the data used in this report}
\small\begin{tabularx}{\textwidth}{@{}K{4.3cm}L>{\raggedright\arraybackslash}p{2.9cm}@{}}
\toprule
\thd{Source} & \thd{Content} & \thd{Coverage} \\
\midrule
NSE combined open interest files & MWPL, open interest, futures-equivalent open interest (from Oct 2025) & Feb 2010 to Oct 2026 \\ \rowrule
NSE securities-in-ban lists & Official daily list of stocks in the ban & Feb 2010 to Oct 2026 \\ \rowrule
NSE F\&O and cash bhavcopies, delivery files & Prices, traded value and quantity, contracts, open interest, delivery & 2015 to Oct 2026 \\ \rowrule
NSE corporate actions & Dividends, splits, bonuses and other events with ex-dates & 2014 to 2026 \\ \rowrule
NSE SLB bhavcopies & Daily stock-lending trades & 2015 to Oct 2026 \\ \rowrule
OECD via FRED & Three-month interbank rate (monthly) & 2011 to 2026 \\
\bottomrule
\end{tabularx}
\exsource{Source: National Stock Exchange of India archives; OECD Main Economic Indicators via FRED.}
\end{table}

\paragraph{Measuring the basis.} For each stock and day I compute the basis \eqref{eq:basis} from the volume-weighted
average prices of the future and of the stock, so that both prices come from trades over the same day:
\begin{equation}\label{eq:basis-data}
m_{it}=\ln\bar F_{it}-\ln\!\big(\bar S_{it}-\mathrm{PV}_{it}(D)\big)-r_t\tau_{it},
\end{equation}
where $\tau_{it}$ is the time to expiry in years and $r_t$ the three-month interbank rate. I use the nearest contract
and switch to the next one when three or fewer trading days remain, require at least five contracts traded, and drop
contracts with a dividend or other corporate event during their life (where the adjustment is uncertain) and the first
20 trading days after a stock enters the F\&O segment. To limit the influence of data errors, the basis is winsorised at
the 1st and 99th percentiles within each year. The final sample has 486,040 stock-days for 371 stocks, 8,851 of them
in a ban.\footnote{The data were checked against the exchange's own rules before use. The ``no fresh positions'' flags
in the open-interest files match the next day's official ban list for 11,753 of 11,764 flags; utilisation was above
95\% on the day before 99.4\% of ban spells began and at or below 80\% on the last day of 99.7\%; and the futures
volume-weighted price lies within the day's high and low for 99.8\% of contracts. \Cref{app:data} gives more detail.}

\subsection{What the data show}
\paragraph{Bans.} From 2015 to October 2026 there were 1,788 ban spells, covering 9,416 stock-days and 141 stocks
(\cref{tab:bans}). Bans are concentrated in a few stocks and come in waves (\cref{fig:banmonth}). A typical spell lasted
four trading days before the October 2025 reform; the 44 spells that began after it have a median of nine.

\begin{table}[t]
\caption{Bans are frequent, concentrated in a few stocks, and longer under the new rule}\label{tab:bans}
\exsub{F\&O stocks, banned stocks, ban spells and ban stock-days by year}
\centering{\footnotesize\setlength{\tabcolsep}{4pt}\begin{tabular}{lrrrrrr}
\toprule
Year & \shortstack[r]{F\&O stocks\\(avg per day)} & \shortstack[r]{Stocks\\banned} & \shortstack[r]{Ban\\spells} & \shortstack[r]{Ban\\stock-days} & \shortstack[r]{Median spell\\(days)} & \shortstack[r]{Longest spell\\(days)} \\
\midrule
2015 & 154 & 16 & 69 & 367 & 4 & 21 \\
2016 & 172 & 27 & 136 & 645 & 4 & 17 \\
2017 & 200 & 43 & 287 & 1586 & 5 & 24 \\
2018 & 206 & 38 & 162 & 810 & 4 & 31 \\
2019 & 176 & 23 & 86 & 508 & 4 & 17 \\
2020 & 140 & 40 & 178 & 756 & 3 & 13 \\
2021 & 163 & 26 & 162 & 789 & 4 & 16 \\
2022 & 196 & 23 & 100 & 457 & 4 & 13 \\
2023 & 187 & 41 & 193 & 942 & 4 & 17 \\
2024 & 185 & 49 & 289 & 1656 & 5 & 22 \\
2025 & 216 & 37 & 93 & 558 & 5 & 16 \\
2026$^{a}$ & 209 & 10 & 33 & 342 & 10 & 29 \\
\midrule
2015 to 2026 & 183 & 141 & 1788 & 9416 & 4 & 31 \\
\bottomrule
\end{tabular}
}
\exsource{$^{a}$ January to 27 October 2026. Spells are counted in the year they start. Source: NSE securities-in-ban lists and F\&O bhavcopies; author's calculations.}
\end{table}

\begin{figure}[t]
\caption{Bans come in waves and became rarer under the new rule}\label{fig:banmonth}
\exsub{Number of stock-days in the ban per month, February 2010 to October 2026}
\centering\includegraphics[width=\linewidth]{fig_ban_month.pdf}
\exsource{Source: NSE securities-in-ban lists; author's calculations.}
\end{figure}

\paragraph{The basis.} The median basis is $+1.1$\bp, so on a typical day futures are priced very close to their
cost-of-carry value. The distribution is skewed: the 10th percentile is $-39.1$\bp\ and the 90th $+26.8$\bp\
(\cref{fig:basisdist}). Cheap futures are thus more common and further from zero than rich ones, which is what one
expects when the trade that would correct cheap futures, short selling the stock, is the costlier one. The basis also
narrows as expiry approaches (\cref{fig:basisexp}), as the carrying part of the costs shrinks with the time left.

\begin{figure}[t]
\caption{The basis is centred near zero with a long lower tail}\label{fig:basisdist}
\exsub{Basis of the analysis sample, 2015 to 2026, in basis points (3.1\% of stock-days lie beyond 100\bp\ from zero and are not drawn)}
\centering\includegraphics[width=\linewidth]{fig_basis_dist.pdf}
\exsource{Source: NSE F\&O and cash bhavcopies, corporate actions, three-month interbank rate; author's calculations.}
\end{figure}

\begin{figure}[t]
\caption{The basis narrows as expiry approaches}\label{fig:basisexp}
\exsub{Median absolute basis of the near-month contract by trading days to expiry (in the last three days the sample uses the next contract)}
\centering\includegraphics[width=0.86\linewidth]{fig_basis_expiry.pdf}
\exsource{Source: as for \cref{fig:basisdist}.}
\end{figure}
